% Challenge dossier: VI.03.C3
\subsection*{Challenge VI.03.C3 --- The union of all coordinate axes}
\label{chall:vi03-c3}
\paragraph{Challenge.}
Let \(X\subseteq\mathbb A_k^n\) be the union of the \(n\) coordinate axes.  Assume \(k\) is
algebraically closed.  Prove
\[
I(X)=(x_ix_j:1\le i<j\le n)
\]
and show
\[
k[X]\cong
\set{(f_1,\dots,f_n)\in k[t]^n:f_1(0)=\cdots=f_n(0)}.
\]
\ifdefined\FullProblemDossiers
\paragraph{Orientation.} A polynomial survives on an axis only through its pure powers in the corresponding variable and its constant term.
\paragraph{Guided hints.} Modulo the pairwise products \(x_ix_j\), every polynomial is congruent to \(c+\sum_i x_i g_i(x_i)\).
\paragraph{Solution.}
Let \(J=(x_ix_j:i<j)\).  Its generators vanish on every coordinate axis, so \(J\subseteq I(X)\).
Modulo \(J\), every mixed monomial vanishes, leaving a unique form
\[
c+\sum_{i=1}^n x_i g_i(x_i).
\]
If this vanishes on every axis, then for each \(i\), the one-variable polynomial
\(c+x_i g_i(x_i)\) vanishes on all \(k\).  Since \(k\) is infinite, it is zero; hence
\(c=0\) and all \(g_i=0\).  Thus \(I(X)=J\).  Restriction to the axes maps a class to a
tuple \((f_1,\dots,f_n)\) with common value at the origin.  Conversely, from such a tuple
with common constant \(c\), the polynomial
\[
c+\sum_i(f_i(x_i)-c)
\]
restricts to the desired tuple.  This proves the ring isomorphism.
\paragraph{Extension.} Identify explicit zero divisors and idempotents, and compare with a disjoint union of \(n\) affine lines.
\paragraph{Provenance.} New challenge generalizing Problem VI.03.P5.
\fi
